% RACER-GT preregistration protocol template (English).
\documentclass[12pt,a4paper,oneside]{article}
% English setup. Compiled with XeLaTeX so that both languages share one engine
% and one font family, which keeps the two PDFs typographically comparable.
\usepackage{fontspec}
\setmainfont{Noto Serif}
\setsansfont{Noto Sans}
\setmonofont{Noto Sans Mono}

% Author-year citations. Loaded here rather than in _common because it must
% precede hyperref, and because the Chinese manuscript uses numeric citations
% that natbib's author-year mode would reject. Bibliographies are hand-written
% thebibliography blocks: biblatex/biber is absent from a BasicTeX install, and
% a document set that only builds on a full TeX Live is not reproducible enough
% for this project.
\usepackage[authoryear,round]{natbib}

% Single source of truth for version and date across every RACER-GT document.
% Regenerate nothing by hand: bump here, rebuild, and every PDF agrees.
\newcommand{\racerversion}{1.1.0}
\newcommand{\racerdate}{2026-07-27}
\newcommand{\racerauthor}{Yi-Hao Lai}
\newcommand{\raceraffil}{Department of Finance, Da-Yeh University}
\newcommand{\raceremail}{yhlai@mail.dyu.edu.tw}
\newcommand{\racerrepo}{https://github.com/yhlai0911/RACER-GT}

% Shared package set and mathematical macros for every RACER-GT document,
% in both languages. Language-specific font and theorem-name setup lives in
% _preamble-zh.tex and _preamble-en.tex.
%
% Font policy: the build depends only on the five Noto families below. They are
% installable on Linux CI (fonts-noto-core, fonts-noto-cjk) and on macOS via
% Homebrew casks. Do not introduce a font outside this set without adding it to
% .github/workflows/docs.yml as well, or the PDF build stops being reproducible.

\usepackage[margin=2.35cm,headheight=15pt]{geometry}
\usepackage{amsmath,amssymb,amsthm,mathtools,bm}
\usepackage{booktabs,longtable,tabularx,array,multirow}
\usepackage{graphicx,float,caption,subcaption}
\usepackage{enumitem}
\usepackage{xcolor}
\usepackage{microtype}
\usepackage{setspace}
\usepackage{fancyhdr}
\usepackage{hyperref}
\usepackage[nameinlink,noabbrev]{cleveref}
\usepackage{algorithm}
\usepackage{algpseudocode}
\usepackage{listings}
\usepackage{tikz}
\usetikzlibrary{arrows.meta,positioning,shapes.geometric,fit,calc}

\hypersetup{
  colorlinks=true,
  linkcolor=black,
  citecolor=black,
  urlcolor=blue,
  pdfauthor={\racerauthor}
}

\pagestyle{fancy}
\fancyhf{}
\rhead{v\racerversion}
\cfoot{\thepage}
\setstretch{1.28}
\setlength{\parskip}{0.35em}
\setlist[itemize]{leftmargin=2em,itemsep=0.25em,topsep=0.3em}
\setlist[enumerate]{leftmargin=2.2em,itemsep=0.25em,topsep=0.3em}

% ---------------------------------------------------------------- mathematics
\newcommand{\E}{\mathbb{E}}
\newcommand{\Var}{\operatorname{Var}}
\newcommand{\Cov}{\operatorname{Cov}}
\newcommand{\tr}{\operatorname{tr}}
\newcommand{\diag}{\operatorname{diag}}
\newcommand{\argmin}{\operatorname*{arg\,min}}
\newcommand{\argmax}{\operatorname*{arg\,max}}
\newcommand{\one}{\bm{1}}
\newcommand{\Rset}{\mathbb{R}}
\newcommand{\R}{\mathbb{R}}
\newcommand{\plim}{\operatorname*{plim}}
\newcommand{\RACER}{\textsc{Racer-GT}}
\newcommand{\indic}[1]{\mathbb{I}\!\left\{#1\right\}}
\newcommand{\Sig}{\bm{\Sigma}}
\newcommand{\wvec}{\bm{w}}
\newcommand{\evec}{\bm{e}}
\newcommand{\Zvec}{\bm{Z}}

% Design facets, used identically in both languages so that a reader can move
% between the Chinese and English documents without re-learning notation.
\newcommand{\facetT}{\mathcal{T}}   % historical date
\newcommand{\facetD}{\mathcal{D}}   % collection day
\newcommand{\facetS}{\mathcal{S}}   % collection stream
\newcommand{\facetR}{\mathcal{R}}   % technical replicate

\lstset{
  basicstyle=\ttfamily\footnotesize,
  breaklines=true,
  frame=single,
  rulecolor=\color{gray!40},
  backgroundcolor=\color{gray!5},
  columns=fullflexible,
  keepspaces=true,
  showstringspaces=false
}


\setlength{\parindent}{1.5em}

\newtheorem{assumption}{Assumption}[section]
\newtheorem{proposition}{Proposition}[section]
\newtheorem{corollary}{Corollary}[section]
\newtheorem{remark}{Remark}[section]
\newtheorem{definition}{Definition}[section]
\newtheorem{lemma}{Lemma}[section]

\lhead{RACER-GT: Protocol and Preregistration}

\title{\bfseries RACER-GT Protocol and Preregistration Template\\[0.4em]
\large What to Lock Before Any Data Are Seen}
\author{\racerauthor\\\small \raceraffil}
\date{\racerdate\ \textbullet\ version \racerversion}

\begin{document}
\maketitle

\begin{abstract}
\noindent This template is meant to be filled in, timestamped, and deposited
\emph{before} collection begins. Its purpose is narrow and specific: to make the
retrieval-side researcher degrees of freedom visible. A Google Trends series
constructed after the analyst has seen the financial results it will be used to
explain is not a measurement; it is a choice. Everything below exists to make
that distinction checkable by someone else.
\end{abstract}

\tableofcontents
\newpage

\section{Study identification}

\begin{itemize}
\item Title, authors, affiliations, contact.
\item Deposit location and timestamp (OSF, AEA RCT Registry, institutional
repository, or a signed Zenodo deposit).
\item RACER-GT version used.
\item \textbf{Protocol hash} produced by \texttt{racergt design}. Record it here.
Any later change to a locked setting changes this hash, which is what makes
drift detectable rather than deniable.
\end{itemize}

\section{Construct and query specification}

State the economic construct first and the query second, in that order, because
the reverse order invites choosing whichever query gives the desired result.

\begin{itemize}
\item The construct the series is intended to measure.
\item Why this keyword or Topic measures it, and what else it plausibly picks up.
\item \texttt{keyword}, \texttt{topic\_or\_term}, \texttt{geo},
\texttt{category}, \texttt{search\_property}, \texttt{language}.
\item Justification for the category choice. Restricting to Finance changes the
construct; so does not restricting. Say which you chose and why.
\item Historical start and end dates, and the baseline window for normalization.
\end{itemize}

\section{Collection design}

\begin{itemize}
\item Collection-day offsets, streams, technical replicates, time slots.
\item What each stream holds fixed: device, browser profile, cookie state,
account state, network route, IP.
\item \textbf{Declaration:} an estimated stream effect will not be reported as
an IP effect unless a crossover factorial design varying IP alone was run.
\item Chunk window length, step, and minimum overlap.
\item The collection route (official API, manual export, institution-approved
client) and its compliance basis.
\end{itemize}

\section{Analysis settings, locked in advance}

Every field of the configuration is part of the protocol. State at minimum:
calibration threshold $c_{\min}$, Huber constant, reference-chunk strategy,
normalization mode, duplicate thresholds, covariance estimator, weight cap,
G-study transformations, benchmark mode and weights.

\section{Acceptance thresholds}

Fill in every threshold \emph{before} collection, and commit to reporting the
resulting status whatever it is.

\begin{table}[H]
\centering\small
\begin{tabularx}{\textwidth}{lXr}
\toprule
Rule & What it protects against & Your value\\
\midrule
\texttt{min\_unique\_pulls} & a batch that is mostly duplicates & \\
\texttt{min\_spectral\_effective\_pulls} & dependence counted as replication & \\
\texttt{max\_zero\_share} & daily frequency unsupportable for a rare term & \\
\texttt{min\_level\_generalizability} & levels that do not reproduce & \\
\texttt{min\_innovation\_generalizability} & daily changes that do not reproduce & \\
\texttt{min\_detection\_kappa} & unstable zero/non-zero pattern & \\
\texttt{max\_component\_share} & one dependence cluster dominating & \\
\texttt{max\_convergence\_mae\_100} & consensus not yet converged & \\
\texttt{max\_benchmark\_standardized\_rmse} & daily and weekly series disagreeing & \\
\bottomrule
\end{tabularx}
\end{table}

\section{Pilot separation}

State explicitly whether an initial monitoring window will be used, and commit
to the following rule: if monitoring leads to \emph{any} change in thresholds,
stitching rules, query, code, or the decision tree, the monitored data become
pilot data and the confirmatory batch restarts from a new Day 0. Without this
commitment the same data design and validate the acceptance rule, which produces
acceptance bias that no amount of subsequent care removes.

\section{Downstream use}

\begin{itemize}
\item The financial model the index will enter, specified now.
\item Whether it enters in levels or differences. If in differences, the
innovation generalizability threshold is the binding one, not the level one.
\item How measurement error will be handled: reliability-corrected OLS, SIMEX,
or both.
\item \textbf{Declaration:} GT processing rules are locked before any return,
volatility, spread, or VaR result is examined.
\end{itemize}

\section{Deviations}

Reserve a section, to be completed after the fact, recording every departure
from this protocol, when it was decided, and why. A protocol with an honest
deviations section is more credible than one with none.

\end{document}
